%===============================================================================
% LaTeX sjabloon voor de bachelorproef toegepaste informatica aan HOGENT
% Meer info op https://github.com/HoGentTIN/latex-hogent-report
%===============================================================================

\documentclass[dutch,dit,thesis]{hogentreport}

%% Pictures to include in the text can be put in the graphics/ folder
\graphicspath{{../graphics/}}

%% For source code highlighting, requires pygments to be installed
%% Compile with the -shell-escape flag!
%% \usepackage[chapter]{minted}
%% If you compile with the make_thesis.{bat,sh} script, use the following
%% import instead:
\usepackage[chapter]{minted}
\usemintedstyle{solarized-light}

%% Formatting for minted environments.
\setminted{%
    autogobble,
    frame=lines,
    breaklines,
    linenos,
    tabsize=4
}

%% Ensure the list of listings is in the table of contents
\renewcommand\listoflistingscaption{%
    \IfLanguageName{dutch}{Lijst van codefragmenten}{List of listings}
}
\renewcommand\listingscaption{%
    \IfLanguageName{dutch}{Codefragment}{Listing}
}
\renewcommand*\listoflistings{%
    \cleardoublepage\phantomsection\addcontentsline{toc}{chapter}{\listoflistingscaption}%
    \listof{listing}{\listoflistingscaption}%
}

% Other packages not already included can be imported here

%%---------- Document metadata -------------------------------------------------
\author{Manu Devloo}
\supervisor{Dhr. T. Clauwaert}
\cosupervisor{Dhr. L. Singier}
\title[Dynamisch Docker-omgevingbeheer voor CTF-events]%
    {\texorpdfstring{Dynamisch Docker-omgevingbeheer\\voor CTF-events\\op on-prem infrastructuur}{Dynamisch Docker-omgevingbeheer voor CTF-events op on-prem infrastructuur}}
\academicyear{2025--2026}
\examperiod{1}
\degreesought{\IfLanguageName{dutch}{Professionele bachelor in de toegepaste informatica}{Bachelor of applied computer science}}
\partialthesis{false} %% To display 'in partial fulfilment'
%\institution{Internshipcompany BVBA.}

%% Add global exceptions to the hyphenation here
\hyphenation{back-slash}

%% The bibliography (style and settings are  found in hogentthesis.cls)
\addbibresource{bachproef.bib}            %% Bibliography file
\addbibresource{../voorstel/voorstel.bib} %% Bibliography research proposal
\defbibheading{bibempty}{}

%% Prevent empty pages for right-handed chapter starts in twoside mode
\renewcommand{\cleardoublepage}{\clearpage}

\renewcommand{\arraystretch}{1.2}
\setlength{\emergencystretch}{2em}

%% Content starts here.
\begin{document}

%---------- Front matter -------------------------------------------------------

\frontmatter

\hypersetup{pageanchor=false} %% Disable page numbering references
%% Render a Dutch outer title page if the main language is English
\IfLanguageName{english}{%
    %% If necessary, information can be changed here
    \degreesought{Professionele Bachelor toegepaste informatica}%
    \begin{otherlanguage}{dutch}%
       \maketitle%
    \end{otherlanguage}%
}{}

%% Generates title page content
\maketitle
\hypersetup{pageanchor=true}

%%=============================================================================
%% Voorwoord
%%=============================================================================

\chapter*{\IfLanguageName{dutch}{Woord vooraf}{Preface}}%
\label{ch:voorwoord}

%% TODO:
%% Het voorwoord is het enige deel van de bachelorproef waar je vanuit je
%% eigen standpunt (``ik-vorm'') mag schrijven. Je kan hier bv. motiveren
%% waarom jij het onderwerp wil bespreken.
%% Vergeet ook niet te bedanken wie je geholpen/gesteund/... heeft

Deze bachelorproef werd uitgevoerd in het kader van de opleiding Professionele bachelor toegepaste informatica aan HOGENT. Mijn interesse in cybersecurity, naast mijn dagelijkse bezigheden binnen softwareontwikkeling, heeft ertoe geleid dat ik dit onderwerp heb gekozen. Binnen dit onderzoek werd gezocht naar een aanpak die zowel veilig als praktisch inzetbaar is op eigen hardware.

Graag bedank ik mijn promotor, dhr.\ Thomas Clauwaert, en mijn co-promotor, dhr.\ Laurens Singier, voor hun begeleiding, feedback en kritische vragen tijdens dit traject.

%%=============================================================================
%% Samenvatting
%%=============================================================================

%%---------- Nederlandse samenvatting -----------------------------------------
%
\IfLanguageName{english}{%
\selectlanguage{dutch}
\chapter*{Samenvatting}
\selectlanguage{english}
}{}

%%---------- Samenvatting -----------------------------------------------------
% De samenvatting in de hoofdtaal van het document

\chapter*{\IfLanguageName{dutch}{Samenvatting}{Abstract}}

Capture The Flag (CTF)-events zijn cybersecuritywedstrijden waarbij deelnemers technische uitdagingen oplossen om punten te behalen. Een deel van die uitdagingen vereist een netwerktoegankelijke service, bijvoorbeeld een webapplicatie of kwetsbare containeromgeving. Op on-premise infrastructuur worden dergelijke challenge-omgevingen vaak nog manueel of met losse scripts beheerd. Dat verhoogt de operationele belasting voor organisatoren en vergroot het risico op misconfiguraties, onvoldoende isolatie tussen deelnemers en resource-uitputting op de host.

Deze bachelorproef onderzoekt hoe Docker-containers die gekoppeld zijn aan CTFd-challenges dynamisch kunnen worden beheerd, aangemaakt en beëindigd op basis van gebruikersacties en configureerbare parameters binnen een volledig self-hosted context. De centrale onderzoeksvraag luidt hoe dit veilig, reproduceerbaar en schaalbaar kan gebeuren zonder afhankelijkheid van betaalde cloudcomponenten of Kubernetes als vereiste.

De methodologie combineert een literatuurstudie met praktijkgericht en iteratief onderzoek. Eerst werden CTFd, Docker, de Docker Engine API en direct toepasbare containerbeveiligingsmaatregelen bestudeerd. Daarna werd een verkennende standalone proof-of-concept gebouwd om container lifecycle management, timeout cleanup, netwerkisolatie, resource-limieten en idempotent startgedrag los van CTFd te valideren. Op basis van die bevindingen werd vervolgens een eigen CTFd-plugin ontwikkeld die een nieuw challenge-type toevoegt en de containerlevenscyclus rechtstreeks koppelt aan acties van spelers en beheerders.

De uiteindelijke plugin-PoC ondersteunt challengeconfiguratie met image, een of meerdere containerpoorten, CPU- en geheugenlimieten, timeout en een maximum aantal gelijktijdige instanties. Daarnaast ondersteunt de plugin het uploaden van Docker-archieven, het automatisch herimporteren van die archieven wanneer een image lokaal ontbreekt, het hergebruiken van een bestaande instantie voor dezelfde speler of hetzelfde team, automatische cleanup bij solve of timeout en een administratieve runtime-pagina binnen CTFd. Voor elke runtime worden basismaatregelen toegepast zoals een read-only root filesystem, capability dropping, \texttt{no-new-privileges}, PID-limieten, CPU- en geheugenlimieten en een afzonderlijk bridge-netwerk per instantie.

De validatie gebeurde in meerdere stappen. De standalone PoC werd gevalideerd met \texttt{pytest}, een Docker-gebaseerde smoke test, isolatie-experimenten, timeoutcontrole en een load snapshot met Locust. De plugin-PoC werd daarna end-to-end getest in same-host modus, in een lokale split-host simulatie en in een reële opstelling met twee afzonderlijke Lima-VM's. Die validatie toont aan dat een CTFd-plugin zowel een lokale Docker-daemon als een remote Docker-host kan aansturen, terwijl spelers correcte runtime-URL's ontvangen voor de afzonderlijke challenge-host.

De conclusie van dit onderzoek is dat een plugin-gebaseerde aanpak binnen CTFd een haalbare en bruikbare oplossingsrichting vormt voor dynamisch containerbeheer bij CTF-events op eigen infrastructuur. Tegelijk blijkt dat verdere productierijpe uitwerking nog aandacht vraagt voor TLS-beveiligde remote Docker-communicatie, uitgebreidere foutinjectie, reverse-proxy-gebaseerde routing en schaalverdeling over meerdere runtime-hosts.


%---------- Inhoud, lijst figuren, ... -----------------------------------------

\tableofcontents

% In a list of figures, the complete caption will be included. To prevent this,
% ALWAYS add a short description in the caption!
%
%  \caption[short description]{elaborate description}
%
% If you do, only the short description will be used in the list of figures

\listoffigures

% If you included tables and/or source code listings, uncomment the appropriate
% lines.
\listoftables

\listoflistings

% Als je een lijst van afkortingen of termen wil toevoegen, dan hoort die
% hier thuis. Gebruik bijvoorbeeld de ``glossaries'' package.
% https://www.overleaf.com/learn/latex/Glossaries

%---------- Kern ---------------------------------------------------------------

\mainmatter{}

% De eerste hoofdstukken van een bachelorproef zijn meestal een inleiding op
% het onderwerp, literatuurstudie en verantwoording methodologie.
% Aarzel niet om een meer beschrijvende titel aan deze hoofdstukken te geven of
% om bijvoorbeeld de inleiding en/of stand van zaken over meerdere hoofdstukken
% te verspreiden!

%%=============================================================================
%% Inleiding
%%=============================================================================

\chapter{\IfLanguageName{dutch}{Inleiding}{Introduction}}%
\label{ch:inleiding}

Capture The Flag (CTF)-events zijn cybersecuritywedstrijden waarbij deelnemers technische uitdagingen oplossen om punten te scoren~\autocite{Innovirtuoso2025}. Ze worden georganiseerd door hogescholen, universiteiten, bedrijven en onafhankelijke gemeenschappen als educatieve en competitieve oefeningen. Elke deelnemer of elk team werkt aan zogenaamde \textit{challenges}: geïsoleerde omgevingen met een kwetsbaarheid die uitgebuit moet worden om een \textit{flag} te bemachtigen.

Naarmate CTF-events populairder worden, groeit ook de behoefte aan robuuste, schaalbare infrastructuur om deze challenges te hosten. Vandaag de dag worden challenges op eigen infrastructuur dikwijls manueel of semi-automatisch beheerd, wat leidt tot operationele overhead, beveiligingsrisico's en schaalbaarheidsknelpunten~\autocite{Singier2025}.

\section{\IfLanguageName{dutch}{Probleemstelling}{Problem Statement}}%
\label{sec:probleemstelling}

Bij het organiseren van CTF-events op eigen (on-premise) infrastructuur is het dynamisch en veilig beheren van containers voor challenges een bijzonder complexe opgave. Organisatoren moeten voor elke challenge een geïsoleerde containeromgeving opzetten, beheren en weer afbreken op basis van gebruikersacties. Dit proces verloopt vandaag grotendeels manueel of met eenvoudige scripts, wat de volgende problemen met zich meebrengt:

\begin{itemize}
  \item Manuele interventie bij het opstarten en stopzetten van containers kost tijd en is foutgevoelig.
  \item Er is geen gegarandeerde isolatie tussen containers van verschillende teams of gebruikers.
  \item Resource-uitputting (CPU, RAM) door ongecontroleerde containers vormt een risico voor de stabiliteit van de host.
  \item Bestaande cloud-gebaseerde platforms (bijv. hosted CTFd) zijn niet altijd geschikt voor organisaties die werken met privédata of op eigen hardware willen draaien.
\end{itemize}

De doelgroep van dit onderzoek zijn CTF-organisatoren en systeembeheerders die events hosten op eigen on-premise infrastructuur en op zoek zijn naar een geautomatiseerde, veilige en schaalbare oplossing.

\section{\IfLanguageName{dutch}{Onderzoeksvraag}{Research question}}%
\label{sec:onderzoeksvraag}

De centrale onderzoeksvraag luidt:

\begin{quote}
  Hoe kunnen Docker-containers die gelinkt zijn aan CTFd-challenges dynamisch worden beheerd, aangemaakt en beëindigd op basis van gebruikersacties en ingestelde parameters, op on-premise infrastructuur?
\end{quote}

Hieruit volgen de volgende deelvragen:

\begin{itemize}
  \item Welke beveiligingsmaatregelen zijn minimaal vereist om container-escapes en misbruik te voorkomen?
  \item Op welke manier kunnen resource-limieten (CPU, RAM) en netwerkisolatie technisch worden afgedwongen?
  \item Hoe kan een CTFd-plugin de Docker Engine API aanspreken zonder zware externe afhankelijkheden?
  \item Hoe valideert men de stabiliteit en veiligheid van de oplossing onder zware belasting?
\end{itemize}

\section{\IfLanguageName{dutch}{Onderzoeksdoelstelling}{Research objective}}%
\label{sec:onderzoeksdoelstelling}

Het beoogde resultaat van deze bachelorproef is een werkende proof-of-concept (PoC) die aantoont hoe CTFd-challenges dynamisch en veilig beheerd kunnen worden op eigen infrastructuur. De PoC omvat een CTFd-plugin die communiceert met de Docker Engine API, resource-limieten afdwingt en netwerkisolatie garandeert. De oplossing is volledig self-hosted, open-source en modulair uitbreidbaar.

\section{\IfLanguageName{dutch}{Scope en afbakening}{Scope}}%
\label{sec:scope}

Dit onderzoek beperkt zich tot on-premise, self-hosted omgevingen. Kubernetes wordt bewust niet als primaire oplossing gekozen omwille van de operationele complexiteit die niet in verhouding staat met de schaal van een typisch CTF-event. Cloudplatforms en betaalde diensten vallen buiten scope. De focus ligt op open-source tooling die draait op standaard Linux-servers.

\section{\IfLanguageName{dutch}{Opzet van deze bachelorproef}{Structure of this bachelor thesis}}%
\label{sec:opzet-bachelorproef}

De rest van deze bachelorproef is als volgt opgebouwd:

In Hoofdstuk~\ref{ch:stand-van-zaken} wordt de achtergrond en het begrippenkader geschetst: wat zijn CTF-events, hoe werkt CTFd, wat zijn de relevante Docker-concepten en welke beveiligingsaspecten spelen een rol?

In Hoofdstuk~\ref{ch:methodologie} wordt de onderzoeksmethodologie toegelicht.

In Hoofdstuk~\ref{ch:huidige-situatie} wordt de huidige situatie (AS-IS) geanalyseerd: hoe worden CTF-events vandaag beheerd en wat zijn de tekortkomingen?

In Hoofdstuk~\ref{ch:gewenste-situatie} wordt de gewenste situatie (TO-BE) beschreven, inclusief de architectuurvisie en de vereisten voor de oplossing.

In Hoofdstuk~\ref{ch:poc} wordt de proof-of-concept uitgewerkt: architectuur, implementatie, beveiliging en schaalbaarheid.

In Hoofdstuk~\ref{ch:validatie} worden de testresultaten besproken: functionele tests, beveiligingstests en load tests.

In Hoofdstuk~\ref{ch:conclusie}, tenslotte, wordt de conclusie gegeven en een antwoord geformuleerd op de onderzoeksvragen, samen met aanbevelingen voor toekomstig werk.

\chapter{\IfLanguageName{dutch}{Achtergrond en begrippenkader}{Background and terminology}}%
\label{ch:stand-van-zaken}

Dit hoofdstuk biedt alle achtergrondkennis die nodig is om de rest van de bachelorproef te begrijpen. Iemand die niet vertrouwd is met CTF-events, Docker of containerbeveiliging, beschikt na het lezen van dit hoofdstuk over een solide basis. Achtereenvolgens komen CTF-events, het CTFd-platform, Docker en containerisatie, containerbeveiliging en load balancing zonder Kubernetes aan bod.

%%----------------------------------------------------------------------
\section{Capture The Flag (CTF)}%
\label{sec:ctf}

Een Capture The Flag (CTF)-event is een cybersecuritywedstrijd waarbij deelnemers technische uitdagingen moeten oplossen om een geheime reeks tekens, de zogenaamde \textit{flag}, te bemachtigen en in te dienen voor punten~\autocite{Innovirtuoso2025}. CTF-events worden ingezet als educatief middel, als rekruteringstool en als competitieve sport binnen de cybersecuritygemeenschap.

\subsection{Types challenges}

CTF-challenges zijn doorgaans opgedeeld in categorieën:

\begin{itemize}
  \item \textbf{Web}: kwetsbaarheden in webapplicaties zoals SQL-injectie, Cross-Site Scripting (XSS) en onveilige autorisatie.
  \item \textbf{Pwn (Binary Exploitation)}: het exploiteren van kwetsbaarheden in gecompileerde binaire bestanden, bijv.\ buffer overflows.
  \item \textbf{Reverse Engineering}: analyse en begrip van gecompileerde software zonder broncode.
  \item \textbf{Cryptografie}: het kraken of omzeilen van cryptografische systemen en protocollen.
  \item \textbf{Forensics}: het reconstrueren van gebeurtenissen uit digitale artefacten zoals logbestanden, geheugenafbeeldingen of netwerkopnamen.
\end{itemize}

\subsection{Deployment van challenges: statisch vs.\ dynamisch}

Challenges kunnen op twee manieren worden uitgerold. Bij een \textbf{statische deployment} wordt één gedeelde instantie van de challenge opgezet die alle deelnemers tegelijk gebruiken. Dit is eenvoudig te beheren, maar brengt risico's mee: één deelnemer kan de omgeving voor alle anderen beïnvloeden. Bij een \textbf{dynamische deployment} krijgt elk team of elke gebruiker een eigen, geïsoleerde containerinstantie. Dit verhoogt de integriteit van de wedstrijd maar vereist automatisering~\autocite{CTFd}.

\subsection{Bekende CTF-platformen}

De meest gebruikte platformen voor het organiseren van CTF-events zijn:

\begin{itemize}
  \item \textbf{CTFd}: het meest wijdverspreide open-source platform, met ondersteuning voor plugins~\autocite{CTFdGitHub}.
  \item \textbf{rCTF}: een alternatief platform ontwikkeld door het redpwn-team~\autocite{rCTF}.
  \item \textbf{kCTF}: een door Google ontwikkeld Kubernetes-gebaseerd platform, ontworpen voor grote-schaal events~\autocite{kCTF}.
\end{itemize}

%%----------------------------------------------------------------------
\section{CTFd als platform}%
\label{sec:ctfd}

CTFd is een open-source CTF-platform, gebouwd in Python met het Flask-framework~\autocite{CTFdGitHub}. Het biedt een gebruiksvriendelijke interface voor organisatoren en deelnemers. Daarnaast is het modulair uitbreidbaar via een plugin-systeem.

\subsection{Architectuur van CTFd}

CTFd bestaat uit een Flask-backend, een relationele database (standaard SQLite, in productie MySQL of PostgreSQL) en een frontend op basis van Bootstrap. De applicatie draait doorgaans achter een reverse proxy zoals Nginx. Alle functionaliteit is bereikbaar via een REST API.

\subsection{Het plugin-systeem}

CTFd ondersteunt plugins die extra functionaliteit toevoegen aan het platform. Een plugin is een Python-pakket dat in de \texttt{CTFd/plugins/}-map geplaatst wordt. Plugins kunnen nieuwe challenge-types registreren, routes toevoegen en bestaande logica uitbreiden~\autocite{CTFdGitHub}.

\subsection{Challenge-beheer en granulariteit}

CTFd ondersteunt standaard statische challenges. Voor dynamische instanties, waarbij elke gebruiker of elk team een eigen containerinstantie krijgt, zijn plugins vereist, zoals de CTFd-Whale plugin~\autocite{CTFdWhale}. De granulariteit (container per gebruiker of per team) is een ontwerpbeslissing met impact op resource-gebruik en isolatie.

%%----------------------------------------------------------------------
\section{Docker en containerisatie}%
\label{sec:docker}

\subsection{Wat is Docker?}

Docker is een platform voor het bouwen, distribueren en draaien van applicaties in containers~\autocite{Arxiv2024}. Een container is een geïsoleerd, lichtgewicht proces dat gebruikmaakt van Linux-kernelmechanismen (namespaces, cgroups) om te worden afgeschermd van andere processen en van de host.

\textbf{Containers versus virtuele machines}: waar een virtuele machine een volledig besturingssysteem emuleert, deelt een container de kernel van de host. Dit maakt containers veel sneller op te starten en minder resource-intensief, maar vergt zorgvuldigere beveiligingsconfiguratie.

\subsection{Docker Engine API}

De Docker Engine API is een RESTful HTTP-API die het mogelijk maakt om programmatisch te communiceren met de Docker-daemon~\autocite{DockerEngineAPI}. Via deze API kunnen containers worden aangemaakt, gestart, gestopt en verwijderd, en kunnen netwerken en volumes worden beheerd. De API is de basis voor alle Docker-tooling, inclusief de Docker CLI.

\subsection{Docker SDK voor Python}

De Docker SDK voor Python is een officiële client-bibliotheek die de Docker Engine API omhult en in Python aanroept~\autocite{DockerSDKPython}. Het is de aanbevolen manier om Docker programmatisch te besturen vanuit Python-applicaties, zoals een CTFd-plugin.

\subsection{Docker Compose}

Docker Compose is een tool voor het definiëren en draaien van multi-container applicaties via een YAML-configuratiebestand~\autocite{DockerCompose}. Voor challenge-configuraties biedt Compose een declaratieve manier om containerparameters (image, netwerk, resource-limieten) te definiëren.

\subsection{Resource management}

Docker biedt mechanismen voor het beperken van resource-gebruik per container, gebaseerd op Linux cgroups:

\begin{itemize}
  \item CPU-limieten: \texttt{--cpus} en \texttt{--cpu-shares}
  \item Geheugenlimieten: \texttt{--memory} en \texttt{--memory-swap}
  \item I/O-limieten via \texttt{blkio}-cgroup-parameters
\end{itemize}

\subsection{Netwerkisolatie}

Docker biedt verschillende netwerkmodi:

\begin{itemize}
  \item \textbf{Bridge}: de standaardmodus waarbij containers communiceren via een virtuele brug, maar gescheiden zijn van de host.
  \item \textbf{None}: geen netwerktoegang.
  \item \textbf{Custom overlay/bridge networks}: containers kunnen worden ingedeeld in geïsoleerde virtuele netwerken.
\end{itemize}

%%----------------------------------------------------------------------
\section{Beveiliging van containers}%
\label{sec:containerbeveiliging}

\subsection{Beveiligingsrisico's bij dynamisch aanmaken van containers}

Wanneer containers dynamisch worden aangemaakt op basis van gebruikersinvoer, ontstaan extra risico's: een kwaadaardige deelnemer kan proberen de configuratie te manipuleren om bevoorrechte toegang te krijgen of om resources te monopoliseren.

\subsection{Direct toepasbare basismaatregelen}

De OWASP Docker Security Cheat Sheet~\autocite{OWASP2024} en artikelen over containerbeveiliging~\autocite{Aikido2025} identificeren een reeks basismaatregelen:

\begin{itemize}
  \item \textbf{Privilegebeheer}: containers draaien standaard als root; dit wordt bij voorkeur vermeden via \texttt{--user}.
  \item \textbf{Read-only bestandssysteem}: via \texttt{--read-only} worden schrijfacties beperkt.
  \item \textbf{Capabilities beperken}: Linux capabilities die niet nodig zijn (bijv.\ \texttt{NET\_ADMIN}, \texttt{SYS\_ADMIN}) worden verwijderd via \texttt{--cap-drop ALL}.
  \item \textbf{Seccomp-profielen}: beperken de systeemoproepen die een container mag uitvoeren.
  \item \textbf{AppArmor/SELinux}: MAC-profielen die extra toegangscontrole bieden.
  \item \textbf{No-new-privileges}: de optie \texttt{--security-opt no-new-privileges} voorkomt privilege-escalatie.
\end{itemize}

Het CIS Docker Benchmark~\autocite{CISDocker} biedt een uitgebreide checklist van aanbevolen beveiligingsinstellingen voor Docker-hosts en -containers.

\subsection{Container escape}

Een container escape is een aanval waarbij een kwaadaardige actor de isolatielaag van de container doorbreekt en toegang verkrijgt tot de host of andere containers. Bekende technieken maken gebruik van kwetsbaarheden in de Docker-daemon, misbruik van gemounte sockets (\texttt{/var/run/docker.sock}) of misconfiguraties (bijv.\ privileged mode)~\autocite{NVD}.

\subsection{Resource exhaustion en DoS-risico's}

Zonder resource-limieten kan een deelnemer de hostserver overbelasten door een excessief aantal containers te starten of door een container alle CPU- en geheugencapaciteit te laten consumeren. Dit vormt een \textit{Denial-of-Service} (DoS)-risico voor alle overige deelnemers.

%%----------------------------------------------------------------------
\section{Load balancing zonder Kubernetes}%
\label{sec:loadbalancing}

\subsection{Waarom geen Kubernetes?}

Kubernetes is het de-facto standaardplatform voor containerorkestratie op schaal~\autocite{kCTF}. Het biedt ingebouwde load balancing, automatische schaling en zelfherstel. Voor kleinschalige CTF-events wegen de voordelen echter niet op tegen de operationele complexiteit: Kubernetes vereist een substantieel cluster, diepgaande kennis en aanzienlijke overhead voor beheer.

\subsection{Alternatieven}

\begin{itemize}
  \item \textbf{Docker Swarm}: Docker's ingebouwde orkestratielaag voor meerdere nodes~\autocite{DockerSwarm}. Eenvoudiger dan Kubernetes, maar minder uitgebreid.
  \item \textbf{Traefik}: een dynamische reverse proxy die automatisch routes configureert op basis van Docker-labels~\autocite{Traefik}. Ideaal als \textit{entry point} voor dynamisch aangemaakte containers.
  \item \textbf{HAProxy}: een betrouwbare, performante TCP/HTTP load balancer~\autocite{HAProxy} voor geavanceerdere routeringsconfiguraties.
  \item \textbf{Nginx}: een veelgebruikte webserver die ook als reverse proxy en load balancer kan fungeren.
\end{itemize}

\subsection{Wanneer is schaling nodig?}

Bij CTF-events varieert de belasting sterk: er zijn piekbelastingen bij de start van een event wanneer veel deelnemers tegelijk challenges starten. De infrastructuur moet deze piekbelasting kunnen opvangen zonder de stabiliteit van het platform te compromitteren.

%%=============================================================================
%% Methodologie
%%=============================================================================

\chapter{\IfLanguageName{dutch}{Methodologie}{Methodology}}%
\label{ch:methodologie}

Dit onderzoek combineert een literatuurstudie met praktijkgericht werk. Het verloopt in vier opeenvolgende fasen.

\textbf{Fase 1 -- Literatuurstudie en begrippenkader} (Hoofdstuk~\ref{ch:stand-van-zaken}): Via literatuurstudie worden de relevante concepten en de state-of-the-art in kaart gebracht. De focus ligt op CTFd, de Docker Engine API, containerbeveiliging en load balancing op on-premise infrastructuur.

\textbf{Fase 2 -- Analyse van de huidige en gewenste situatie} (Hoofdstukken~\ref{ch:huidige-situatie} en~\ref{ch:gewenste-situatie}): De huidige aanpak voor het beheren van CTF-challenges op eigen infrastructuur wordt geanalyseerd (AS-IS). Op basis hiervan en de literatuurstudie worden functionele en niet-functionele vereisten opgesteld en een architectuurvisie uitgetekend (TO-BE).

\textbf{Fase 3 -- Proof-of-Concept} (Hoofdstuk~\ref{ch:poc}): Er wordt een werkende PoC ontwikkeld: een CTFd-plugin die via de Docker Engine API containers dynamisch beheert. De PoC implementeert resource-limieten, netwerkisolatie en de basisbeveiligingsmaatregelen die in de literatuurstudie zijn geïdentificeerd. De schaalbaarheid wordt getest met een multi-VM opzet en een reverse proxy.

\textbf{Fase 4 -- Validatie en testing} (Hoofdstuk~\ref{ch:validatie}): De PoC wordt gevalideerd aan de hand van een testplan met functionele tests, beveiligingstests (container escape pogingen, resource exhaustion) en load tests met tools zoals k6~\autocite{k6} of Locust~\autocite{Locust}.


% Corpus van de bachelorproef
%%=============================================================================
%% Huidige situatie (AS-IS)
%%=============================================================================

\chapter{\IfLanguageName{dutch}{Huidige situatie}{Current situation}}%
\label{ch:huidige-situatie}

Dit hoofdstuk analyseert hoe CTF-events vandaag worden georganiseerd op eigen infrastructuur. Het brengt de problemen en beperkingen van de huidige aanpak in kaart als vertrekpunt voor de gewenste oplossing.

%%----------------------------------------------------------------------
\section{Organisatie van CTF-events op eigen infrastructuur}%
\label{sec:asis-organisatie}

% TODO: Beschrijf hoe CTF-events vandaag typisch worden georganiseerd op
% on-premise infrastructuur. Welke tools en processen worden gebruikt?
% Verwijs naar gesprekken met organisatoren of bestaande documentatie.
% Bronverwijzingen: \autocite{CTFd}, \autocite{Singier2025}

%%----------------------------------------------------------------------
\section{Manuele container deployments}%
\label{sec:asis-manueel}

% TODO: Bespreek hoe challenges vandaag worden uitgerold:
% - handmatige docker run / docker-compose commando's
% - scripts zonder foutafhandeling
% - geen geautomatiseerd opruimen van gestopte containers
% Welke problemen en beperkingen zijn hier aan verbonden?

%%----------------------------------------------------------------------
\section{Overzicht van bestaande tools en hun tekortkomingen}%
\label{sec:asis-tools}

% TODO: Geef een overzicht van de tools die vandaag gebruikt worden en
% beschrijf hun beperkingen:
% - CTFd zonder plugin: enkel statische challenges
% - CTFd-Whale: \autocite{CTFdWhale} - bestaande plugin, maar beperkingen?
% - Handmatige scripts: niet schaalbaar, geen isolatiegaranties
% - Hosted platforms (bijv. CTFd cloud): niet self-hosted, kostprijs

%%----------------------------------------------------------------------
\section{Conclusie: wat ontbreekt er?}%
\label{sec:asis-conclusie}

% TODO: Formuleer op basis van de analyse welke elementen ontbreken in de
% huidige aanpak. Dit vormt de directe input voor de TO-BE beschrijving.
% Denk aan: automatisering, isolatie, resource-beheer, schaalbaarheid,
% gebruiksgemak voor organisatoren.

%%=============================================================================
%% Gewenste situatie (TO-BE)
%%=============================================================================

\chapter{\IfLanguageName{dutch}{Gewenste situatie}{Target situation}}%
\label{ch:gewenste-situatie}

Op basis van de tekortkomingen die in het vorige hoofdstuk werden geïdentificeerd, beschrijft dit hoofdstuk de vereisten voor de oplossing en de architectuurvisie voor het gewenste systeem.

%%----------------------------------------------------------------------
\section{Vereisten voor de oplossing}%
\label{sec:tobe-vereisten}

\subsection{Functionele vereisten}

% TODO: Lijst de functionele vereisten op:
% - Automatisch aanmaken van een containerinstantie per gebruiker/team bij
%   het starten van een challenge
% - Automatisch stoppen en verwijderen na afloop of timeout
% - Configureerbare resource-limieten per challenge (CPU, RAM)
% - Netwerkisolatie tussen instanties van verschillende gebruikers
% - Ondersteuning voor meerdere challenge-types
% - Gebruiksvriendelijke interface voor organisatoren om challenges toe te voegen

\subsection{Niet-functionele vereisten}

% TODO: Lijst de niet-functionele vereisten op:
% - Volledig self-hosted en open-source
% - Geen Kubernetes als primaire vereiste
% - Schaalbaar over meerdere VM's
% - Veilig: minimale aanvalsoppervlakte, container escapes voorkomen
% - Beheerbaar: monitoring en logging van actieve instanties
% - Reproduceerbaar: gedocumenteerde installatie-instructies

%%----------------------------------------------------------------------
\section{Architectuurvisie}%
\label{sec:tobe-architectuur}

% TODO: Beschrijf de architectuur van het gewenste systeem:
% - CTFd communiceert via plugin met Docker Engine API
%   (\autocite{DockerEngineAPI}, \autocite{DockerSDKPython})
% - Containers worden automatisch aangemaakt/vernietigd op basis van
%   gebruikersacties in CTFd
% - Resource-limieten worden afgedwongen bij aanmaak
% - Netwerkisolatie via custom Docker bridge networks
% - Reverse proxy (Traefik \autocite{Traefik} of Nginx) als entry point
% - Systeem is modulair: nieuwe challenge-types toevoegen zonder kerncode
%   te wijzigen

% TODO: Voeg hier een componentendiagram in.

%%----------------------------------------------------------------------
\section{Granulariteitsbeslissing: container per gebruiker of per team}%
\label{sec:tobe-granulariteit}

% TODO: Bespreek de voor- en nadelen van beide opties:
% - Per gebruiker: maximale isolatie, meer resource-gebruik
% - Per team: minder instanties, risico op interne interferentie
% Verantwoord de keuze die in de PoC wordt geïmplementeerd.

%%----------------------------------------------------------------------
\section{Schaalbaarheidsvereisten}%
\label{sec:tobe-schaalbaarheid}

% TODO: Beschrijf wat het systeem moet aankunnen:
% - Hoeveel gelijktijdige deelnemers?
% - Hoeveel actieve containers tegelijk?
% - Hoe worden containers verdeeld over meerdere hosts?
% - Docker Swarm \autocite{DockerSwarm} of andere orchestratielaag?

%%=============================================================================
%% Proof-of-Concept
%%=============================================================================

\chapter{\IfLanguageName{dutch}{Proof-of-Concept}{Proof of Concept}}%
\label{ch:poc}

Dit hoofdstuk beschrijft de praktische uitwerking van het onderzoek. De nadruk ligt op de uiteindelijke CTFd-plugin, omdat die rechtstreeks aansluit bij de onderzoeksvraag: challenge-containers dynamisch beheren vanuit CTFd. Voorafgaand daaraan werd echter een beperkte standalone PoC gebouwd om de kernmechanismen sneller te testen. Die verkennende stap wordt kort besproken omdat ze verschillende ontwerpkeuzes voor de plugin heeft onderbouwd.

%%----------------------------------------------------------------------
\section{Architectuuroverzicht}%
\label{sec:poc-architectuur}

De uiteindelijke PoC bestaat uit een CTFd-instantie waarin een eigen plugin een nieuw challenge-type toevoegt. Die plugin spreekt de Docker-daemon aan via de Docker SDK voor Python~\autocite{DockerSDKPython} en beheert de volledige levenscyclus van per-gebruiker of per-team challenge-instanties. Dezelfde plugin kan zowel een lokale Docker-daemon via de Unix-socket aanspreken als een remote Docker API op een andere host. Runtime-metadata, zoals actieve instanties en image-archieven, worden lokaal bijgehouden zodat de plugin ook cleanup en administratieve acties kan uitvoeren.

De opzet bevat vier kernonderdelen:

\begin{itemize}
  \item \textbf{CTFd}: het wedstrijdplatform met authenticatie, challenges en scoreverwerking.
  \item \textbf{CTFd-plugin}: de eigenlijke orkestratielaag die start-, stop- en cleanupacties uitvoert.
  \item \textbf{Docker-daemon}: de runtime waarin challenge-containers en hun netwerken worden aangemaakt; in split-host modus draait die op een afzonderlijke host.
  \item \textbf{Runtime-opslag}: een lichte SQLite-opslag voor actieve instanties, logs en metadata van geuploade image-archieven.
\end{itemize}

Tijdens de implementatie werd eerst vertrokken van een single-host opstelling om de integratie beheersbaar te houden. Nadien is dezelfde architectuur uitgebreid naar een split-host model waarin het Docker-control endpoint en de spelergerichte runtime-host afzonderlijk configureerbaar zijn. Binnen de validatie van deze bachelorproef werd die tweede modus zowel lokaal gesimuleerd als op twee afzonderlijke VM's getest.

\begin{figure}[H]
  \centering
  \resizebox{\textwidth}{!}{%
  \begin{tikzpicture}[
    node distance=1.05cm and 1.25cm,
    font=\small,
    box/.style={draw, rounded corners, align=center, text width=2.75cm, minimum height=1.25cm, fill=black!3},
    plugin/.style={draw, rounded corners, align=center, text width=3.2cm, minimum height=1.7cm, fill=black!6},
    store/.style={draw, rounded corners, align=center, text width=3.2cm, minimum height=1.25cm, fill=white},
    arrow/.style={->, thick},
    note/.style={font=\scriptsize, align=center, fill=white, inner sep=1.5pt}
  ]
    \node[box] (player) {\textbf{Deelnemer of\\ organisator}\\start, stop, solve\\beheeracties};
    \node[box, right=of player] (ctfd) {\textbf{CTFd}\\challenge-flow\\beheerpagina};
    \node[plugin, right=of ctfd] (plugin) {\textbf{Containerized plugin}\\policy- en lifecyclelogica\\plugin-hooks\\views};
    \node[box, right=of plugin] (docker) {\textbf{Docker-daemon}\\containers\\netwerken};
    \node[box, below=of docker] (runtime) {\textbf{Challenge-instantie}\\per gebruiker\\of team};
    \node[store, below=of plugin] (storage) {\textbf{Runtime-opslag}\\SQLite\\image-archieven\\logs en metadata};

    \draw[arrow] (player) -- node[note, above=3pt, pos=0.6] {start / stop / solve} (ctfd);
    \draw[arrow] (ctfd) -- node[note, above=3pt, pos=0.58] {plugin-hooks\\views} (plugin);
    \draw[arrow] (plugin) -- node[note, above=3pt, pos=0.55] {Docker SDK} (docker);
    \draw[arrow] (docker) -- node[note, right=3pt] {run / rm\\network create} (runtime);
    \draw[arrow] (plugin) -- node[note, left=3pt] {metadata\\image-archieven} (storage);
  \end{tikzpicture}
  }
  \caption[Architectuur van de plugin-PoC]{Architectuur van de plugin-PoC. CTFd blijft het centrale platform, terwijl de plugin de lifecycle van challenge-containers en de bijhorende metadata beheert. De Docker-daemon kan lokaal of op een afzonderlijke runtime-host draaien.}
  \label{fig:poc-architectuur}
\end{figure}

%%----------------------------------------------------------------------
\section{Verkennende standalone PoC}%
\label{sec:poc-standalone}

Voor de plugin werd uitgewerkt, werd eerst een kleine standalone orchestrator gebouwd in Flask. Het doel daarvan was niet om een alternatieve eindarchitectuur naar voren te schuiven, maar om drie vragen snel te beantwoorden:

\begin{itemize}
  \item Kan een Python-service containers starten, stoppen en opruimen via de Docker SDK?
  \item Welke basisbeveiligingsmaatregelen zijn praktisch toepasbaar zonder de challenge-flow onnodig complex te maken?
  \item Hoe kunnen lifecycle-regels zoals timeouts, quota en idempotent starten eenvoudig gevalideerd worden?
\end{itemize}

De standalone PoC bood een challenge-registry, een eenvoudige webinterface, API-endpoints voor start en stop, een timeout-reaper en een beperkte logginglaag. Uit die stap volgden een aantal keuzes die rechtstreeks in de plugin zijn meegenomen:

\begin{itemize}
  \item gebruik van de Docker SDK in plaats van shell-calls naar de Docker CLI;
  \item idempotent startgedrag voor dezelfde gebruiker of hetzelfde team;
  \item standaardresourceprofielen per challenge;
  \item cleanup van containers en netwerken bij timeout of expliciete stop;
  \item netwerkisolatie per instantie.
\end{itemize}

De standalone PoC was dus vooral een technische voorstudie die toeliet om de risicovolste runtime-aspecten vroeg te toetsen. Zodra die basis werkte, werd de focus verschoven naar de CTFd-plugin, omdat die nauwer aansluit bij het doel van dit onderzoek: de containerlevenscyclus rechtstreeks koppelen aan acties in het CTF-platform.

\begin{table}[htbp]
  \centering
  \caption[Vergelijking van beide PoC's]{Vergelijking tussen de verkennende standalone PoC en de uiteindelijke CTFd-plugin PoC.}
  \label{tab:poc-vergelijking}
  \begin{tabularx}{\textwidth}{@{}lXX@{}}
    \toprule
    \textbf{Aspect} & \textbf{Standalone PoC} & \textbf{CTFd-plugin PoC} \\
    \midrule
    Hoofddoel & Snelle validatie van lifecycle-mechanismen en beveiligingsmaatregelen. & Integratie van die mechanismen in de echte challenge-flow van CTFd. \\
    Interface & Eigen Flask-dashboard en REST-API. & Challenge-type en beheerpagina binnen CTFd. \\
    Bewijswaarde & Toont dat Docker-orkestratie technisch werkt in isolatie. & Toont dat gebruikersacties in CTFd rechtstreeks runtimegedrag kunnen sturen. \\
    Belangrijkste output & Inzichten over idempotent starten, timeout cleanup en isolatie. & Werkende plugin met uploadflow, solve-cleanup en administratieve controle. \\
    Rol in deze bachelorproef & Verkennende tussenstap om ontwerpkeuzes te onderbouwen. & Hoofdresultaat dat de onderzoeksvraag het dichtst benadert. \\
    \bottomrule
  \end{tabularx}
\end{table}

%%----------------------------------------------------------------------
\section{Implementatie van de CTFd-plugin}%
\label{sec:poc-plugin}

\subsection{Custom challenge-type en configuratie}

De plugin voegt een nieuw challenge-type \texttt{containerized} toe aan CTFd. Een organisator kan zo per challenge runtime-parameters configureren zonder externe scripts of manuele \texttt{docker run}-commando's. De belangrijkste parameters zijn:

\begin{itemize}
  \item Docker-image of geupload image-archief;
  \item een of meerdere containerpoorten;
  \item CPU-limiet;
  \item geheugenlimiet;
  \item timeout;
  \item maximum aantal gelijktijdige instanties.
\end{itemize}

Naast een klassieke image-referentie ondersteunt de plugin ook het uploaden van een \texttt{.tar}, \texttt{.tar.gz} of \texttt{.tgz}-archief uit \texttt{docker save}. Dat archief wordt lokaal opgeslagen, in de Docker-daemon geimporteerd en aan de challenge gekoppeld. Daardoor kan een organisator ook zonder externe registry een image beschikbaar maken voor de PoC. Wanneer de lokaal geimporteerde image later ontbreekt, kan de plugin het gekoppelde archief automatisch opnieuw importeren bij een volgende runtime-start. Die eigenschap bleek bijzonder nuttig in de smoke tests, omdat het daarmee mogelijk werd om de volledige upload- en herimportflow reproduceerbaar te valideren.

\subsection{Start- en stopflow voor deelnemers}

Wanneer een deelnemer een containerized challenge opent, kan die via het challenge-scherm een eigen runtime-instantie starten. De plugin vertaalt die actie naar een Docker-start met de juiste challenge-configuratie. In individuele modus gebeurt dit per gebruiker; in teammodus per team. Daarbij gelden twee belangrijke beleidsregels:

\begin{itemize}
  \item voor dezelfde gebruiker of hetzelfde team wordt een bestaande actieve instantie hergebruikt in plaats van een nieuwe te starten;
  \item zodra het ingestelde maximum aantal instanties is bereikt, weigert de plugin bijkomende starts.
\end{itemize}

Na een succesvolle start krijgt de deelnemer een of meerdere toegangspunten terug met de gepubliceerde hostpoort(en). Die URL's worden opgebouwd op basis van een publiek configureerbare runtime-host, zodat de plugin ook correct blijft werken wanneer CTFd en Docker niet op dezelfde machine draaien. De plugin bewaart tegelijk metadata zoals eigenaar, starttijd, vervaltijd, container-ID, netwerknaam en poortbindingen. Wanneer de deelnemer expliciet stopt, wordt de container verwijderd en wordt ook het bijhorende Docker-netwerk opgeruimd.

\subsection{Timeout cleanup en cleanup bij solve}

De plugin bevat een achtergrondproces dat periodiek controleert welke instanties vervallen zijn. Wanneer een timeout bereikt is, wordt de container geforceerd gestopt en worden de runtime-resources vrijgemaakt. De eindstatus van zo'n instantie wordt als verlopen geregistreerd zodat achteraf zichtbaar blijft waarom ze stopte.

Daarnaast grijpt de plugin ook in op het moment dat een challenge correct opgelost wordt. In dat geval wordt de actieve instantie van de betrokken gebruiker of het betrokken team automatisch stopgezet. Die koppeling is belangrijk voor CTF-contexten: een deelnemer hoeft een oplossing niet eerst manueel te stoppen en de infrastructuur blijft niet onnodig resources verbruiken nadat de flag al werd ingediend.

\begin{figure}[H]
  \centering
  \resizebox{0.95\textwidth}{!}{%
  \begin{tikzpicture}[
    node distance=0.9cm and 0.95cm,
    flowbox/.style={draw, rounded corners, align=center, minimum width=2.65cm, minimum height=0.85cm, fill=black!3},
    arrow/.style={->, thick}
  ]
    \node[flowbox] (open) {Challenge\\ openen};
    \node[flowbox, right=of open] (policy) {Quota en\\ timeout check};
    \node[flowbox, right=of policy] (start) {Container\\ starten};
    \node[flowbox, right=of start] (access) {Runtime-URL\\ tonen};

    \node[flowbox, below=of access] (active) {Actieve\\ instantie};
    \node[flowbox, below left=of active] (stop) {Stop of\\ solve};
    \node[flowbox, below right=of active] (timeout) {Timeout\\ bereikt};
    \node[flowbox, below=1.05cm of active] (cleanup) {Container en\\ netwerk opruimen};

    \draw[arrow] (open) -- (policy);
    \draw[arrow] (policy) -- (start);
    \draw[arrow] (start) -- (access);
    \draw[arrow] (access) -- (active);
    \draw[arrow] (active) -- (stop);
    \draw[arrow] (active) -- (timeout);
    \draw[arrow] (stop) -- (cleanup);
    \draw[arrow] (timeout) -- (cleanup);
  \end{tikzpicture}
  }
  \caption[Start-, stop- en cleanupflow]{Start-, stop- en cleanupflow voor een containerized challenge. Dezelfde cleanupstap wordt gebruikt bij manueel stoppen, een correcte solve en timeout.}
  \label{fig:runtime-flow}
\end{figure}

\subsection{Administratieve runtime-interface}

Voor organisatoren voorziet de plugin een aparte beheerpagina binnen CTFd. Daarop zijn twee soorten informatie zichtbaar:

\begin{itemize}
  \item een overzicht van alle containerized challenges met hun resourceprofiel;
  \item een runtime-overzicht van actieve en historische instanties.
\end{itemize}

Vanuit die interface kunnen beheerders een individuele instantie geforceerd stoppen, alle actieve instanties van een challenge stoppen en de timeout-reaper manueel uitvoeren. Dit verhoogt de beheerbaarheid van de PoC aanzienlijk, omdat de organisator niet rechtstreeks op de Docker-host hoeft in te loggen om runtime-problemen op te lossen.

\clearpage

%%----------------------------------------------------------------------
\section{Beveiliging in de plugin-PoC}%
\label{sec:poc-beveiliging}

De plugin past bij elke containerstart een aantal direct toepasbare basismaatregelen toe, in lijn met de OWASP Docker Security Cheat Sheet~\autocite{OWASP2024} en het CIS Docker Benchmark~\autocite{CISDocker}. Concreet gaat het om de volgende Docker-instellingen:

\begin{itemize}
  \item \textbf{Read-only root filesystem}: het hoofdbestandssysteem van de container wordt alleen-lezen gemount. De challenge kan daardoor niet zomaar systeembestanden, configuratie of applicatiebestanden in de container aanpassen. Dat beperkt de impact van een exploit die probeert persistentie of extra tooling in de container weg te schrijven.
    \item \textbf{\texttt{cap\_drop=['ALL']}}: Linux-capabilities zijn fijnmazige rechten die processen extra kernelmogelijkheden geven, zoals netwerkconfiguratie aanpassen of systeembeheeracties uitvoeren. Door alle capabilities standaard weg te nemen, start de container met zo weinig mogelijk privileges. Alleen wanneer een challenge aantoonbaar extra rechten nodig heeft, zou een organisator daar bewust van kunnen afwijken.
  \item \textbf{\texttt{security\_opt=['no-new-privileges:true']}}: deze instelling verhindert dat processen binnen de container via mechanismen zoals setuid-binaries extra privileges verkrijgen. Zelfs wanneer een programma binnen de container normaal gezien met verhoogde rechten zou starten, mag het proces geen nieuwe privileges krijgen bovenop wat het al had.
  \item \textbf{PID-limiet}: het aantal processen dat binnen de container mag draaien wordt begrensd. Dit voorkomt eenvoudige fork bombs of challenges die per ongeluk zeer veel processen aanmaken en zo de Docker-host belasten.
  \item \textbf{CPU- en geheugenlimieten}: per challenge wordt vastgelegd hoeveel CPU-tijd en RAM een container maximaal mag gebruiken. Daardoor kan één foutieve of kwaadwillige challenge-instantie niet alle resources van de host opeisen en andere deelnemers of CTFd zelf verstoren.
  \item \textbf{Tijdelijke \texttt{tmpfs}-mount voor \texttt{/tmp}}: de map \texttt{/tmp} blijft beschrijfbaar, maar wordt in geheugen geplaatst en verdwijnt wanneer de container stopt. Zo kunnen applicaties die tijdelijke bestanden nodig hebben blijven werken, terwijl permanente writes naar het containerbestandssysteem vermeden worden.
  \item \textbf{Bridge-netwerk per runtime-instantie}: elke challenge-instantie krijgt een eigen Docker-netwerk. Containers van verschillende gebruikers of teams delen daardoor geen standaard netwerksegment en kunnen elkaar niet zomaar via containernaam of intern IP-adres ontdekken.
\end{itemize}

Die combinatie beperkt het risico op privilege-escalatie, laterale beweging en resource-uitputting. De keuze voor per-instantie netwerken zorgt ervoor dat deelnemers niet zomaar andere challenge-instanties kunnen ontdekken of aanspreken. Voor split-host deployments ondersteunt de plugin daarnaast een expliciet configureerbaar gepubliceerd poortbereik, zodat firewallregels op de runtime-host begrensd kunnen blijven. Belangrijk is wel dat deze PoC een basisbeveiliging demonstreert en geen volledig production-ready hardeningprofiel oplevert. Zo is ondersteuning voor aangepaste seccomp-profielen nog geen afgewerkt onderdeel van de implementatie en werd remote Docker-communicatie in de validatie enkel op wegwerpomgevingen zonder TLS getest.

\clearpage

%%----------------------------------------------------------------------
\section{Deployment- en schaalbaarheidsoverwegingen}%
\label{sec:poc-schaalbaarheid}

De huidige plugin-PoC ondersteunt twee deploymentmodi. In same-host modus draait CTFd samen met de plugin op dezelfde host als de Docker-daemon en wordt de lokale Unix-socket gebruikt. In split-host modus blijft CTFd het control plane, maar wordt de Docker-daemon op een andere host aangesproken via een afzonderlijk control endpoint. De spelergerichte runtime-host kan daarbij verschillen van het Docker API endpoint. Binnen die scope toont de PoC aan dat een plugin-gebaseerde aanpak haalbaar is voor:

\begin{itemize}
  \item dynamisch starten en stoppen van challenge-containers;
  \item afdwingen van limieten per challenge;
  \item automatisch cleanupgedrag bij timeout en solve;
  \item administratieve controle vanuit CTFd;
  \item scheiding tussen control plane en runtime-host in een split deployment.
\end{itemize}

De split-host ondersteuning werd in deze bachelorproef ook effectief gevalideerd. Eerst gebeurde dat lokaal via een Compose-opstelling met een aparte \texttt{docker:dind}-service en een smoke-runner binnen hetzelfde netwerk. Daarna volgde een reële validatie op twee aparte Lima-VM's: één VM voor CTFd en de plugin, en één VM voor de challenge-containers. Daarbij werden de gepubliceerde challenge-poorten begrensd tot een vast bereik, zodat de runtime-host eenvoudiger afgeschermd kon worden.

Voor verdere schaalvergroting blijven wel nog een aantal vragen over:

\begin{itemize}
  \item hoe remote Docker-communicatie in productie TLS-beveiligd en operationeel beheersbaar wordt ingericht;
  \item hoe het systeem reageert wanneer de runtime-host tijdelijk onbereikbaar wordt tijdens start, stop of cleanup;
  \item hoe dynamische routing best wordt gecombineerd met een reverse proxy in plaats van louter hostpoortmapping;
  \item hoe meerdere runtime-hosts evenwichtig kunnen worden aangestuurd vanuit dezelfde plugin.
\end{itemize}

%%=============================================================================
%% Validatie & Testing
%%=============================================================================

\chapter{\IfLanguageName{dutch}{Validatie en testing}{Validation and testing}}%
\label{ch:validatie}

Dit hoofdstuk beschrijft het testplan en de testresultaten waarmee de correctheid, veiligheid en schaalbaarheid van de PoC worden aangetoond.

%%----------------------------------------------------------------------
\section{Testplan}%
\label{sec:validatie-testplan}

\subsection{Doelstellingen}

% TODO: Beschrijf wat het testplan beoogt te valideren:
% - Werkt het dynamisch aanmaken en beëindigen van containers correct?
% - Is de isolatie tussen instanties gegarandeerd?
% - Zijn de beveiligingsmaatregelen effectief?
% - Houdt het systeem stand onder zware belasting?

\subsection{Testomgeving}

% TODO: Beschrijf de testomgeving:
% - Hardware/VM-specificaties
% - Betrokken software en versies (CTFd, Docker, plugin, OS)
% - Netwerktopologie van de testopstelling

%%----------------------------------------------------------------------
\section{Functionele tests}%
\label{sec:validatie-functioneel}

\subsection{Aanmaken en beëindigen van containers}

% TODO: Beschrijf de testscenario's:
% - Deelnemer start een challenge: wordt de container correct aangemaakt?
% - Deelnemer stopt de challenge: wordt de container correct verwijderd?
% - Timeout: wordt de container automatisch gestopt na de ingestelde tijd?
% - Meerdere deelnemers starten tegelijk dezelfde challenge.

\subsection{Isolatie tussen gebruikers}

% TODO: Beschrijf hoe wordt geverifieerd dat containers van verschillende
% deelnemers elkaar niet kunnen bereiken:
% - Netwerktest: kan container A container B pingen?
% - Kan een deelnemer de omgeving van een andere deelnemer manipuleren?

%%----------------------------------------------------------------------
\section{Beveiligingstests}%
\label{sec:validatie-beveiliging}

\subsection{Container escape pogingen}

% TODO: Beschrijf welke aanvalstechnieken werden getest:
% - Misbruik van gemounte volumes
% - Privilege escalation via capabilities
% - Docker socket exposure
% Welke beschermingen verhinderden succesvolle escapes?

\subsection{Resource exhaustion simulaties}

% TODO: Beschrijf hoe resource exhaustion werd gesimuleerd:
% - Fork bomb binnen een container
% - Geheugenoverbelasting
% Hoe reageerde het systeem? Werden de resource-limieten afgedwongen?

%%----------------------------------------------------------------------
\section{Load testing}%
\label{sec:validatie-loadtest}

\subsection{Testopzet}

% TODO: Beschrijf de load test:
% - Gebruikte tool: k6 \autocite{k6} of Locust \autocite{Locust}
% - Testscenario: X gelijktijdige gebruikers starten een challenge
% - Duur en ramp-up profiel

\subsection{Meetpunten}

% TODO: Welke metrics werden gemeten?
% - Latency bij het aanmaken van een container
% - CPU- en RAM-gebruik op de host(s)
% - Foutpercentage (mislukte container-starts)
% - Doorvoer (containers per seconde)

\subsection{Resultaten en conclusies}

% TODO: Presenteer de testresultaten in tabellen of grafieken.
% Wat zijn de grenzen van de PoC? Bij welk aantal gelijktijdige gebruikers
% treedt er degradatie op?


%%=============================================================================
%% Conclusie
%%=============================================================================

\chapter{Conclusie}%
\label{ch:conclusie}

Deze bachelorproef onderzocht hoe Docker-containers die gelinkt zijn aan CTFd-challenges dynamisch beheerd, aangemaakt en beëindigd kunnen worden op basis van gebruikersacties en ingestelde parameters op on-premise infrastructuur. Op basis van de literatuurstudie, de verkennende standalone PoC en de uiteindelijke plugin-PoC kan geconcludeerd worden dat een plugin-gebaseerde aanpak binnen CTFd hiervoor een haalbare oplossingsrichting is.

%%----------------------------------------------------------------------
\section{Beantwoording van de onderzoeksvragen}%
\label{sec:conclusie-onderzoeksvragen}

De centrale onderzoeksvraag kan positief beantwoord worden. Het is mogelijk om vanuit CTFd challenge-containers dynamisch te starten en opnieuw te stoppen, zolang de containerlevenscyclus expliciet wordt gekoppeld aan challenge-acties en runtime-beleid. In de uitgewerkte PoC gebeurt dat via een eigen \texttt{containerized} challenge-type dat Docker-containers start met vooraf ingestelde limieten en ze achteraf ook weer opruimt.

Ook de deelvragen kunnen op basis van de resultaten beantwoord worden:

\begin{itemize}
  \item \textbf{Welke beveiligingsmaatregelen zijn minimaal vereist?} Een bruikbaar basisprofiel bestaat uit een read-only root filesystem, het droppen van Linux capabilities, \texttt{no-new-privileges}, CPU- en geheugenlimieten, een PID-limiet en netwerkisolatie per instantie.
  \item \textbf{Hoe worden resource-limieten en netwerkisolatie afgedwongen?} De PoC gebruikt Docker runtime-opties voor CPU, geheugen en PID-limieten en maakt per runtime-instantie een eigen bridge-netwerk aan.
  \item \textbf{Hoe communiceert CTFd met Docker?} In deze bachelorproef gebeurt dat via een eigen plugin in Python die de Docker SDK voor Python gebruikt om containers, netwerken en image-imports aan te sturen. Die communicatie werkt zowel tegen een lokale Docker-socket als tegen een remote Docker API op een andere host.
  \item \textbf{Hoe werd de stabiliteit gevalideerd?} Via een combinatie van geautomatiseerde tests, smoke tests, timeout-validatie, isolatie-experimenten, een beperkte loadtest met Locust en split-host validatie in zowel een lokale simulatie als een reële twee-VM-opstelling.
\end{itemize}

%%----------------------------------------------------------------------
\section{Beperkingen van de PoC}%
\label{sec:conclusie-beperkingen}

Hoewel de plugin-PoC de kern van de onderzoeksvraag beantwoordt, zijn er duidelijke beperkingen. De huidige implementatie werd wel gevalideerd in same-host modus en in een split deployment over twee VM's, maar ze toont nog niet hoe de oplossing zich gedraagt in een bredere meerhost-opstelling met taakverdeling over meerdere runtime-hosts. Ook reverse-proxy-routing, TLS-beveiligde remote Docker-communicatie, geavanceerde production hardening en aangepaste seccomp-profielen maken nog geen afgerond deel uit van de PoC.

Daarnaast was de standalone PoC nuttig als exploratief tussenresultaat, maar ze is niet bedoeld als eindarchitectuur. De standalone stap diende vooral om risico's vroeg te reduceren en technische keuzes te onderbouwen voordat de pluginintegratie werd gebouwd.

%%----------------------------------------------------------------------
\section{Aanbevelingen voor productiegebruik}%
\label{sec:conclusie-aanbevelingen}

Voor productiegebruik volstaat de huidige PoC nog niet zonder bijkomende maatregelen. Organisaties die verder willen bouwen op deze aanpak, moeten minstens rekening houden met:

\begin{itemize}
  \item hardening van de Docker-host en het netwerksegment waarin challenge-containers draaien;
  \item versleutelde en streng afgeschermde communicatie indien Docker niet lokaal op dezelfde host draait;
  \item een expliciet beheer van gepubliceerde poortbereiken en firewallregels op de runtime-host;
  \item bijkomende monitoring en alerting voor mislukte starts, cleanupfouten en capaciteitspieken;
  \item operationele procedures voor imagebeheer, updates en incidentrespons.
\end{itemize}

De plugin-aanpak blijft daarbij een passende basis, omdat ze de gebruikerservaring en het beheer in CTFd centraliseert en tegelijk voldoende controle biedt over runtime-beleid.

%%----------------------------------------------------------------------
\section{Mogelijke uitbreidingen}%
\label{sec:conclusie-uitbreidingen}

De belangrijkste volgende stap is het hardenen van de reeds bewezen split-host architectuur. Daarbij moeten minstens vier vragen beantwoord worden:

\begin{itemize}
  \item hoe remote Docker-communicatie via TLS en passende toegangscontrole concreet wordt opgezet;
  \item of een dunne tussenlaag toch wenselijk is tussen plugin en Docker-host voor extra beveiliging of foutafhandeling;
  \item wat de impact is van die opsplitsing op opstartlatency en cleanupbetrouwbaarheid onder hogere belasting;
  \item hoe het systeem reageert wanneer de Docker-host tijdelijk onbereikbaar is.
\end{itemize}

Andere interessante uitbreidingen zijn een reverse-proxy-gebaseerd toegangspunt, uitgebreidere loadtesting en verdere production hardening.


%---------- Bijlagen -----------------------------------------------------------

\appendix

\chapter{Onderzoeksvoorstel}

Het onderwerp van deze bachelorproef is gebaseerd op een onderzoeksvoorstel dat vooraf werd beoordeeld door de promotor. Dat voorstel is opgenomen in deze bijlage.

\section*{Samenvatting}

Dit onderzoeksvoorstel behandelt het dynamisch beheer van Docker-containers voor Capture The Flag (CTF)-events op on-premise infrastructuur. De probleemstelling richt zich op het aanmaken en beëindigen van containers op basis van gebruikersacties. De hoofdonderzoeksvraag onderzoekt hoe dit proces kan worden geautomatiseerd met behulp van CTFd en Docker. Daarbij wordt gewerkt binnen de randvoorwaarden van een volledig self-hosted omgeving, zonder terugkerende licentiekosten en met een voorkeur voor open-source componenten. Via een literatuurstudie en een proof-of-concept worden de beveiligingsaspecten, resource-limieten en communicatie tussen CTFd en Docker onderzocht. Het doel is een schaalbare oplossing te ontwerpen die misbruik helpt voorkomen en de performantie bewaakt.

% Verwijzing naar het bestand met de inhoud van het onderzoeksvoorstel
\input{../voorstel/voorstel-inhoud}

%%---------- Andere bijlagen --------------------------------------------------

\chapter{Installatie- en beheerinstructies}
\label{app:installatie}

Deze bijlage beschrijft hoe beide PoC's gereproduceerd kunnen worden. De standalone PoC is opgenomen als voorstudie; de plugin-PoC is het hoofdresultaat van de bachelorproef.

\section*{Standalone PoC}

De standalone PoC kan gebruikt worden om de kernmechanismen los van CTFd te reproduceren.

\begin{enumerate}
  \item Zorg dat Docker Desktop of Docker Engine actief is en dat Docker Compose v2 beschikbaar is.
  \item Navigeer naar de map \texttt{POC/}.
  \item Start de end-to-end smoke test via \texttt{./scripts/smoke\_test.sh}.
  \item Open nadien het dashboard op \texttt{http://127.0.0.1:8000}.
\end{enumerate}

Voor lokale tests zijn Python 3.12 en de dependencies uit \texttt{app/requirements-dev.txt} nodig. De tests kunnen uitgevoerd worden met \texttt{app/.venv/bin/pytest -q tests}.

\section*{CTFd-plugin PoC}

De plugin-PoC is de primaire reproductieroute voor deze bachelorproef.

\begin{enumerate}
  \item Zorg dat Docker actief is en dat Docker Compose v2 beschikbaar is.
  \item Navigeer naar de map \texttt{POC-CTFd/}.
  \item Start de same-host validatie via \texttt{./scripts/smoke\_test.sh}.
  \item Start de split-host validatie via \texttt{./scripts/smoke\_test\_split.sh} wanneer de control plane en runtime-host logisch gescheiden moeten worden.
  \item Open CTFd op \texttt{http://127.0.0.1:8001}.
\end{enumerate}

De same-host smoke test initialiseert CTFd, uploadt een Docker-archief, maakt een \texttt{containerized} challenge aan, start runtime-instanties voor meerdere spelers, controleert bereikbaarheid en valideert cleanup na solve, manuele stop en timeout. De split-host smoke test doorloopt dezelfde functionele flow, maar laat CTFd communiceren met een afzonderlijke Docker-runtime.

Een reële twee-VM reproductie is bijkomend gedocumenteerd in \path{POC-CTFd/docs/03_real-vm-validation.md}. Daarin wordt een Lima-opstelling beschreven met CTFd en plugin op een eerste VM en de challenge-containers op een tweede VM.

\section*{Beheer}

Voor de plugin-PoC gebeurt beheer grotendeels vanuit CTFd zelf. De plugin voegt een administratieve runtime-pagina toe waarop actieve instanties zichtbaar zijn. Daar kan een beheerder:

\begin{itemize}
  \item individuele instanties geforceerd stoppen;
  \item alle actieve instanties van een challenge tegelijk stopzetten;
  \item de timeout-reaper manueel uitvoeren;
  \item per challenge het resourceprofiel raadplegen.
\end{itemize}

Voor de standalone PoC bestaan daarnaast helper-scripts, zoals \path{cleanup_runtime.sh}, om alle beheerde runtime-artifacten op te ruimen.

\chapter{Configuratiebestanden}
\label{app:configuratie}

Deze bijlage vat de belangrijkste configuratie-elementen samen die voor de PoC relevant zijn.

\section*{Plugin-configuratie}

De plugin voegt per challenge de volgende configuratievelden toe:

\begin{itemize}
  \item image of geupload Docker-archief;
  \item een of meerdere containerpoorten;
  \item CPU-limiet;
  \item geheugenlimiet;
  \item timeout in seconden;
  \item maximum aantal gelijktijdige instanties.
\end{itemize}

Daarnaast ondersteunt de plugin optioneel een geupload Docker-archief dat lokaal wordt opgeslagen en nadien in de Docker-daemon geimporteerd of opnieuw geimporteerd kan worden.

\section*{Runtime-configuratie via environment variables}

De plugin gebruikt onder meer environment variables voor:

\begin{itemize}
  \item \texttt{CTFD\_CONTAINER\_DOCKER\_HOST} voor het Docker-control endpoint;
  \item \texttt{CTFD\_CONTAINER\_DOCKER\_TLS\_VERIFY} voor TLS-verificatie richting Docker;
  \item \texttt{CTFD\_CONTAINER\_DOCKER\_CERT\_PATH} voor clientcertificaten richting Docker;
  \item \texttt{CTFD\_CONTAINER\_PUBLIC\_HOST} voor spelergerichte runtime-URL's;
  \item \texttt{CTFD\_CONTAINER\_PUBLIC\_SCHEME} voor het URL-schema van die runtime-URL's;
  \item \texttt{CTFD\_CONTAINER\_PUBLISHED\_PORT\_MIN} en \texttt{CTFD\_CONTAINER\_PUBLISHED\_PORT\_MAX} voor een optioneel begrensd poortbereik;
  \item \texttt{CTFD\_CONTAINER\_DB\_PATH} voor de runtime-database;
  \item \texttt{CTFD\_CONTAINER\_IMAGE\_ARCHIVE\_DIR} voor opslag van geuploade image-archieven;
  \item \texttt{CTFD\_CONTAINER\_REAPER\_INTERVAL} voor het interval van de timeout-reaper.
\end{itemize}

\section*{Container hardening}

Bij het starten van een challenge-container worden standaard de volgende runtime-opties toegepast:

\begin{itemize}
  \item read-only root filesystem;
  \item \texttt{cap\_drop=['ALL']};
  \item \texttt{security\_opt=['no-new-privileges:true']};
  \item \texttt{tmpfs} voor \texttt{/tmp};
  \item CPU-, geheugen- en PID-limieten;
  \item een afzonderlijk bridge-netwerk per instantie.
\end{itemize}

\section*{Niet-opgenomen configuratie}

Reverse-proxy-configuratie en taakverdeling over meerdere runtime-hosts zijn bewust niet opgenomen als afgewerkte configuratiebestanden. De split-host opstelling met één aparte runtime-host werd wel gevalideerd, maar productierijpe TLS-configuratie en meerhost-orkestratie blijven toekomstig werk.

\chapter{Testresultaten in detail}
\label{app:testresultaten}

Deze bijlage bundelt de belangrijkste commando's en samengevatte resultaten die in de validatiehoofdstukken gebruikt werden.

\section*{Standalone PoC}

\textbf{Pytest}
\begin{verbatim}
cd POC
app/.venv/bin/pytest -q tests
\end{verbatim}

Resultaat: \texttt{4 passed in 0.60s}

\textbf{Smoke test}
\begin{verbatim}
cd POC
./scripts/smoke_test.sh
\end{verbatim}

Resultaat: de demo-image en orchestrator werden opgebouwd, twee instanties werden gestart, de HTTP-endpoints waren bereikbaar en beide instanties werden nadien correct gestopt.

\textbf{Load snapshot}
\begin{verbatim}
cd POC
app/.venv/bin/locust \
  -f tests/locustfile.py \
  --host=http://127.0.0.1:8000 \
  --headless -u 10 -r 5 -t 10s
\end{verbatim}

Samenvatting: 128 requests, 84 failures, telkens door \texttt{409 CONFLICT} na het bereiken van \texttt{max\_instances}.

\section*{Plugin-PoC}

\textbf{End-to-end smoke test}
\begin{verbatim}
cd POC-CTFd
./scripts/smoke_test.sh
\end{verbatim}

Resultaat: CTFd werd geinitialiseerd, een Docker-archief werd geupload, een \texttt{containerized} challenge werd aangemaakt, de lokaal geimporteerde image werd verwijderd om herimport te forceren, meerdere spelers startten runtimes, \texttt{max\_instances} werd afgedwongen, een tweede start hergebruikte dezelfde instantie, manuele stop en timeout cleanup werden bevestigd en na solve werd de runtime automatisch opgeruimd met stopreden \texttt{challenge-solved}.

\textbf{Split-host smoke test}
\begin{verbatim}
cd POC-CTFd
./scripts/smoke_test_split.sh
\end{verbatim}

Resultaat: dezelfde challenge-flow werkte ook wanneer CTFd met de plugin via een remote Docker endpoint communiceerde. De gegenereerde runtime-URL's verwezen naar de afzonderlijke runtime-host en de gepubliceerde poorten bleven binnen het ingestelde bereik.

\textbf{Re\"ele twee-VM validatie}
\begin{verbatim}
Zie POC-CTFd/docs/03_real-vm-validation.md
\end{verbatim}

Resultaat: op 24 maart 2026 werd bevestigd dat CTFd op een eerste Lima-VM challenge-containers kon starten, stoppen en opruimen op een tweede Lima-VM. Na afloop bleven geen plugin-beheerde containers of netwerken achter op de runtime-VM.

\section*{Interpretatie}

De gedetailleerde outputs tonen vooral aan dat de gekozen lifecycle-regels effectief worden afgedwongen. Ze tonen ook aan dat split deployment technisch haalbaar is met dezelfde pluginarchitectuur. Ze vormen wel nog geen bewijs voor een volledig afgewerkte productiearchitectuur met TLS-beveiligde remote Docker-communicatie, reverse proxy en taakverdeling over meerdere runtime-hosts.

\chapter{Onderzoekslog}
\label{app:onderzoekslog}

Dit appendix documenteert de belangrijkste beslissingen die tijdens het onderzoek werden genomen, met hun motivatie. Het geeft inzicht in het verloop van het onderzoeksproces en de keuzes die de uiteindelijke proof-of-concept hebben gevormd.

%%----------------------------------------------------------------------
\section*{Fase 1 -- Literatuurstudie en begrippenkader}

\textbf{Beslissing: afbakening van de literatuurstudie}\\
Bij de start van het onderzoek werd beslist de literatuurstudie te richten op vier thema's: CTF-events en het CTFd-platform, Docker en containerisatie, containerbeveiliging en load balancing zonder Kubernetes. Dit stelde de lezer in staat de rest van de bachelorproef te volgen zonder voorkennis. Diepgaande theoretische beveiligingsmodellen (zoals formele verificatie of CVE-analyse) werden bewust buiten scope gelaten ten voordele van direct toepasbare best practices~\autocite{OWASP2024,CISDocker}.

\textbf{Beslissing: focus op ``low-hanging fruit'' beveiliging}\\
Na het bestuderen van de OWASP Docker Security Cheat Sheet~\autocite{OWASP2024} en het CIS Docker Benchmark~\autocite{CISDocker} werd gekozen om de focus te leggen op praktisch toepasbare basismaatregelen: capabilities droppen, read-only bestandssysteem, seccomp-profielen, no-new-privileges en resource-limieten. Deze maatregelen bieden een goede balans tussen beveiliging en implementatiecomplexiteit voor een CTF-context.

%%----------------------------------------------------------------------
\section*{Fase 2 -- Analyse van de huidige en gewenste situatie}

\textbf{Beslissing: Kubernetes niet als primaire oplossing}\\
Kubernetes werd vroeg in het onderzoek overwogen als orkestratieplatform. Na analyse werd besloten het niet als primaire oplossing te gebruiken omdat de operationele complexiteit (clusteropzet, beheer, leercurve) niet in verhouding staat met de schaal van een doorsnee CTF-event. Als alternatief werd gefocust op Docker Swarm~\autocite{DockerSwarm} en Traefik~\autocite{Traefik} als lichtere alternatieven. Dit sluit ook aan bij de randvoorwaarde van eenvoudige inzetbaarheid op eigen infrastructuur.

\textbf{Beslissing: granulariteit -- container per gebruiker of per team}\\
Bij de architectuuranalyse werd de keuze tussen een instantie per gebruiker en een instantie per team bestudeerd. Per gebruiker biedt maximale isolatie maar verhoogt het aantal containers snel. Per team beperkt het resource-gebruik maar verlaagt de isolatiegarantie. Uiteindelijk werd gekozen voor een flexibele aanpak: in teammodus wordt één instantie per team voorzien, bij individuele deelname één per gebruiker. Dit koppelt de granulariteit aan het wedstrijdmodel en vermijdt onnodige containerexplosie.

\textbf{Beslissing: Docker SDK voor Python vs.\ directe API-calls}\\
Voor de integratie tussen de CTFd-plugin en de Docker-daemon werden twee opties afgewogen: de officiële Docker SDK voor Python~\autocite{DockerSDKPython} en directe HTTP-aanroepen naar de Docker Engine API~\autocite{DockerEngineAPI}. De SDK biedt een hogere abstractie, minder foutgevoeligheid en is de aanbevolen aanpak voor Python-applicaties. Directe API-calls zouden meer controle geven maar ook meer boilerplate en foutafhandeling vergen. Gekozen werd voor de SDK, tenzij er specifieke beperkingen opduiken tijdens de implementatie.

\textbf{Beslissing: reverse proxy niet in de eerste implementatie}\\
Een reverse proxy met dynamische routing, bijvoorbeeld via Traefik~\autocite{Traefik}, bleef architecturaal interessant, maar werd niet in de eerste implementatiefase opgenomen. De reden daarvoor was pragmatisch: eerst moest duidelijk worden dat container lifecycle management, cleanup en basisbeveiliging betrouwbaar werkten. Reverse-proxy-routing werd daarom verschoven naar een latere uitbreidingsfase.

\textbf{Beslissing: eerst een verkennende standalone stap, daarna plugin-integratie}\\
Hoewel de einddoelstelling altijd plugin-gebaseerd bleef, werd tijdens de analyse beslist eerst een beperkte standalone PoC te bouwen. Die keuze liet toe om de kern van container lifecycle management los van CTFd te valideren en technische risico's vroeg te reduceren. Zodra die basis stabiel genoeg bleek, werd de focus verschoven naar de eigenlijke CTFd-plugin.

\textbf{Beslissing: bestaande plugin CTFd-Whale als referentie, niet als basis}\\
De bestaande CTFd-Whale plugin~\autocite{CTFdWhale} werd bestudeerd als referentie voor de architectuuranalyse. Er werd besloten geen fork of uitbreiding te maken, maar een eigen implementatie te schrijven. CTFd-Whale voldoet niet volledig aan de randvoorwaarden van dit onderzoek (volledig self-hosted, open-source, controleerbare beveiligingsinstellingen) en een eigen implementatie biedt meer controle over de beveiligings- en resourceconfiguratie.

%%----------------------------------------------------------------------
\section*{Fase 3 -- Uitwerking en herori\"entatie}

\textbf{Beslissing: standalone PoC behouden als voorstudie, niet als eindarchitectuur}\\
Na de eerste experimenten bleek dat de standalone PoC bruikbaar was voor het valideren van timeouts, netwerkisolatie, logging en basisresourcebeleid. Tegelijk werd duidelijk dat deze opzet minder goed aansloot bij de centrale onderzoeksvraag, omdat de gebruikersinteractie dan buiten CTFd bleef. De standalone PoC werd daarom behouden als onderzoeksstap en niet als finale oplossingsrichting.

\textbf{Beslissing: plugin behouden als voorkeursrichting}\\
Na de realisatie van de plugin-PoC werd beslist om de plugin als voorkeursrichting naar voren te schuiven. De plugin koppelt gebruikersacties rechtstreeks aan challenge-runtimes, houdt de beheerinterface binnen CTFd en sluit daardoor goed aan bij de context van CTF-events op eigen infrastructuur.

\section*{Fase 4 -- Validatie en afronding}

\textbf{Beslissing: split deployment als aparte validatiestap uitvoeren}\\
Tijdens de eerste uitwerking werd split deployment nog niet meteen als resultaat gepresenteerd. Er werd beslist om dit niet speculatief in de PoC-hoofdstukken op te nemen, maar pas na een afzonderlijke validatie. In een latere iteratie werd daarvoor zowel een lokale split-host simulatie als een re\"ele twee-VM test op Lima uitgevoerd. Op die manier kon de haalbaarheid van de architectuur alsnog onderbouwd worden met concrete meet- en testresultaten.

\textbf{Beslissing: remote Docker-communicatie aantonen, maar niet als production-ready voorstellen}\\
De extra validatie toonde aan dat de plugin een remote Docker-host kan aansturen en correcte spelergerichte runtime-URL's kan genereren. Tegelijk werd beslist dit niet als volledig productieklare oplossing te presenteren, omdat de gebruikte testopstellingen nog geen TLS-beveiligde remote API, uitgebreide foutinjectie of taakverdeling over meerdere runtime-hosts omvatten.

%%---------- Backmatter, referentielijst ---------------------------------------

\backmatter{}

\setlength\bibitemsep{2pt} %% Add Some space between the bibliograpy entries
\printbibliography[heading=bibintoc]

\end{document}
